\chapter{Modèles probabilistes et incertitude}\label{ch:ml:probabilistic-models-and-uncertainty}

Deux prévisions annoncent toutes deux qu'une action gagnera cinq points de base demain. L'une s'accompagne d'un écart-type de 80
points de base, l'autre avec 250~: le même rendement attendu un jour calme et un jour agité. Un livre qui dimensionne les deux
positions de la même manière détient trois fois le risque par unité de rendement attendu dans le second cas, et un rapport de risque qui n'affiche que
la prévision ponctuelle ne permet pas de les distinguer. Toute prévision d'un rendement est une distribution. Ce chapitre entraîne des modèles
qui le disent, les notent au moyen de pertes qui récompensent les distributions honnêtes, construisent des intervalles dont la couverture résiste à un
changement de régime, et mesure ce que vaut la variance lorsque les positions sont dimensionnées.

\section{Distributions prédictives et règles de score propres}

\begin{definition}[Distribution prédictive, règle de score propre]\label{def:ml:probabilistic-models-and-uncertainty:scoring}
Une \emph{distribution prédictive}\index{distribution prédictive} (predictive distribution) est une prévision de toute la distribution conditionnelle
$F(\cdot\mid x)$ de la cible. Une règle de score $S(F, y)$ attribue une perte à une prévision $F$ lorsque $y$ se réalise~; c'est une
\emph{règle de score propre}\index{règle de score propre} (proper scoring rule) si son espérance sous la vraie distribution $G$ est minimisée
par la prévision $F = G$, de sorte que l'honnêteté soit optimale (Gneiting et Raftery, 2007). La log-vraisemblance négative est propre~;
de même que les deux règles ci-dessous.
\end{definition}

\begin{definition}[Perte pinball, score de probabilité classé continu]\label{def:ml:probabilistic-models-and-uncertainty:crps}
La \emph{perte pinball}\index{perte pinball} (pinball loss) d'une prévision $q$ du quantile d'ordre $\alpha$ est
$\ell_\alpha(y, q) = \max\{\alpha(y - q), (\alpha - 1)(y - q)\}$. Le
\emph{score de probabilité classé continu}\index{score de probabilité classé continu} (continuous ranked probability score) d'une distribution de prévision $F$ est
$\mathrm{CRPS}(F, y) = \int(F(z) - \mathbf 1\{z\ge y\})^2\,dz = 2\int_0^1\ell_\alpha(y, F^{-1}(\alpha))\,d\alpha$, exprimé dans les
unités de la cible~; pour une prévision gaussienne, elle admet une forme close.
\end{definition}

\begin{proposition}[La perte pinball élicite le quantile]\label{prop:ml:uncert:pinball}
Si $Y$ admet une fonction de répartition continue $G$, $\E\,\ell_\alpha(Y, q)$ est minimisée en $q = G^{-1}(\alpha)$.
\end{proposition}

\begin{proof}
$\frac{d}{dq}\E\,\ell_\alpha(Y, q) = -\alpha\P(Y > q) + (1 - \alpha)\P(Y < q) = G(q) - \alpha$, qui s'annule en
$G^{-1}(\alpha)$, négative avant et positive après.
\end{proof}

Les données du chapitre sont les séries quotidiennes de \texttt{firm.mlsynth}~: vingt actifs, dix ans, volatilité GARCH, chocs de Student-$t$
à cinq degrés de liberté, une dérive injectée, et un régime à partir du jour 2\,100 dans lequel chaque choc est 1.6 fois
plus grand. La cible est le rendement du lendemain~; les variables sont la lecture de la dérive, les rendements sur 5 et 20 jours, le rapport
entre la volatilité réalisée courte et la volatilité réalisée longue, et une volatilité EWMA, toutes connues à la clôture. Jours d'entraînement 60--1\,511, jours
de calibration 1\,512--1\,763, jours de test 1\,764--2\,518~: le régime change à peu près à mi-parcours de la période de test.

\section{Pertes quantiles et distributionnelles}

\begin{definition}[Régression quantile, réseau de densité de mélange]\label{def:ml:probabilistic-models-and-uncertainty:quantile}
La \emph{régression quantile}\index{régression quantile} (quantile regression) estime le quantile conditionnel d'ordre $\alpha$ en minimisant la moyenne
de la \omterm{def:ml:probabilistic-models-and-uncertainty:crps}{perte pinball} $\mathcal L_n = \frac1n\sum_i\ell_\alpha(y_i, f_\theta(x_i))$ (Koenker et Bassett, 1978)~; avec plusieurs
niveaux, elle décrit toute la distribution. Un \emph{réseau de densité de mélange}\index{réseau de densité de mélange} (mixture density network) produit les
poids, moyennes et écarts-types d'un mélange de $K$ gaussiennes, et il est entraîné par leur log-vraisemblance négative
(Bishop, 1994).
\end{definition}

Trois modèles distributionnels sont ajustés~: des arbres à \omterm{def:ml:trees-and-boosting:boosting}{gradient boosting} pour dix-neuf niveaux de quantile (5~\% à 95~\%), un réseau
dont la sortie comporte une moyenne et une log-variance, entraîné par la log-vraisemblance gaussienne négative (cinq graines, moyennées en \omterm{def:ml:neural-networks-for-noisy-tabular-data:ensemble}{ensemble profond},
chapitre~7), et un \omterm{def:ml:probabilistic-models-and-uncertainty:quantile}{réseau de densité de mélange} à deux composantes. Le \Cref{tab:ml:uncert:scores} les évalue avec le CRPS calculé sur les
dix-neuf quantiles et la couverture de leurs intervalles à 90~\%, avant et après le changement de régime, à côté de
l'approximation gaussienne de la vraie distribution conditionnelle.

\begin{table}[htbp]
\centering\small
\setlength{\tabcolsep}{4pt}%
\begin{tabular}{@{}>{\raggedright\arraybackslash}p{4.4cm}>{\centering\arraybackslash}p{2.3cm}>{\centering\arraybackslash}p{2.3cm}>{\centering\arraybackslash}p{2.3cm}>{\centering\arraybackslash}p{2.3cm}@{}}
\toprule
 & \multicolumn{2}{c}{avant le jour 2\,100} & \multicolumn{2}{c}{après le jour 2\,100}\\
\cmidrule(lr){2-3}\cmidrule(l){4-5}
modèle & CRPS (pb) & couverture à 90~\% & CRPS (pb) & couverture à 90~\%\\
\midrule
quantiles boostés & 78.9 & 89.4~\% & 136.3 & 85.0~\%\\
réseau gaussien (une graine) & 79.0 & 90.6~\% & 135.3 & 87.0~\%\\
ensemble de cinq & 79.1 & 91.1~\% & 135.4 & 87.6~\%\\
\omterm{def:ml:probabilistic-models-and-uncertainty:quantile}{réseau de densité de mélange} & 78.8 & 89.7~\% & 136.1 & 86.1~\%\\
\midrule
vraie distribution (approximation gaussienne) & 78.9 & 91.2~\% & 135.5 & 91.3~\%\\
\bottomrule
\end{tabular}
\caption{Prévisions distributionnelles du rendement du lendemain sur les jours de test de vingt actifs synthétiques. CRPS calculé sur
dix-neuf quantiles (5~\% à 95~\%), en points de base de rendement. Données~: \texttt{ml\_uncert.scores}.}\label{tab:ml:uncert:scores}
\end{table}

Avant la rupture, les quatre modèles font aussi bien que l'approximation de la vraie distribution, à moins d'un quart de point de base
près du CRPS~: la volatilité qu'ils doivent prévoir est persistante et la variable EWMA en porte l'essentiel. Après la rupture, le CRPS
de tous quasiment double, comme celui de la vraie distribution, mais leurs intervalles couvrent 85--88~\% des jours au lieu de 90~\%~: les
variables s'adaptent avec un retard, et les modèles ont été ajustés sur un monde plus calme. Les prévisions de moyenne sont une autre affaire~; un
rapport signal sur bruit de quelques centièmes ne déplace guère aucun de ces scores, qui sont dominés par la volatilité.

\section{Deux types d'incertitude}

\begin{definition}[Incertitude aléatoire et incertitude épistémique]\label{def:ml:probabilistic-models-and-uncertainty:aleatoric}
L'\emph{incertitude aléatoire}\index{incertitude aléatoire} (aleatoric uncertainty) est le caractère aléatoire de la cible conditionnellement aux entrées, que davantage
de données ne peuvent supprimer~; l'\emph{incertitude épistémique}\index{incertitude épistémique} (epistemic uncertainty) est l'incertitude du modèle sur la
fonction elle-même, que davantage de données réduisent. Dans un \omterm{def:ml:neural-networks-for-noisy-tabular-data:ensemble}{ensemble profond}, la moyenne des variances prédites par les membres estime
la première, et la variance de leurs moyennes la seconde~; la variance prédictive est leur somme.
\end{definition}

Sur les jours de test du chapitre, la part épistémique représente 0.09~\% de la variance prédictive. Les rendements quotidiens sont presque entièrement
aléatoires~: rien de ce qu'un modèle apprend ne les resserrera beaucoup. La part épistémique compte là où le modèle extrapole
(entrées hors de la plage d'entraînement, un nouvel instrument) et dans les situations à peu d'étiquettes du chapitre~10~; c'est la partie
qu'un ensemble, et non un réseau isolé, peut signaler (Kendall et Gal, 2017).

\section{Prédiction conforme}

\begin{definition}[Prédiction conforme]\label{def:ml:probabilistic-models-and-uncertainty:conformal}
La \emph{prédiction conforme}\index{prédiction conforme} (conformal prediction) transforme toute prévision en ensembles de prédiction munis d'une garantie
de couverture~: un score $s_i = s(x_i, y_i)$ (par exemple $|y_i - \hat y_i|$, ou ce même score divisé par un
écart-type prédit), est calculé sur $n$ exemples de calibration, et l'ensemble pour un nouveau $x$ est $\{y: s(x, y)\le\hat q\}$ avec $\hat q$ le
$\lceil(n+1)(1 - \alpha)\rceil$-ième plus petit score de calibration (prédiction conforme scindée~; Vovk, Gammerman et Shafer, 2005). La
version adaptative met à jour le niveau en ligne, $\alpha_{t+1} = \alpha_t + \gamma(\alpha - \mathrm{err}_t)$, afin de conserver le
taux d'erreur à long terme à $\alpha$ lorsque les données dérivent (Gibbs et Candès, 2021).
\end{definition}

\begin{theorem}[Couverture de la prédiction conforme scindée]\label{thm:ml:uncert:conformal}
Si les scores de calibration et le nouveau score sont échangeables, $\P(s_{n+1}\le\hat q)\ge1 - \alpha$.
\end{theorem}

\begin{proof}
Par échangeabilité, le rang de $s_{n+1}$ parmi les $n + 1$ scores est uniforme sur $\{1,\dots,n+1\}$ (les égalités étant départagées au hasard).
$s_{n+1}\le\hat q$ dès que son rang vaut au plus $k = \lceil(n+1)(1-\alpha)\rceil$, ce qui a pour probabilité
$k/(n+1)\ge1 - \alpha$.
\end{proof}

L'échangeabilité est précisément ce que brise un changement de régime (\cref{fig:ml:uncert:coverage}). Autour de la moyenne
prédictive des arbres boostés, la \omterm{def:ml:probabilistic-models-and-uncertainty:conformal}{prédiction conforme} scindée sur l'année de calibration couvre 90.7~\% des jours de test avant la rupture et 74.0~\% après celle-ci,
à largeur fixe. Diviser le score par l'écart-type prédit de l'ensemble permet à la largeur de suivre la volatilité~:
90.3~\% avant, 86.9~\% après. La \omterm{def:ml:probabilistic-models-and-uncertainty:conformal}{prédiction conforme} adaptative, appliquée aux mêmes scores normalisés et mise à jour jour par jour pendant la
période de test avec $\gamma = 0.005$, couvre 89.9~\% avant et 90.1~\% après (\cref{fig:ml:uncert:band})~: elle renonce à la
garantie à échantillon fini au profit d'une garantie de long terme qui résiste à la dérive.

\begin{omfigure}
\begin{tikzpicture}
\begin{axis}[width=12cm,height=5.4cm,ybar,bar width=7pt,ylabel={couverture de l'intervalle à 90~\% (\%)},
  symbolic x coords={boosted quantiles,ensemble of 5,mixture density,conformal raw,conformal normalised,conformal adaptive},
  xtick=data,xticklabels={{quantiles\\boostés},{ensemble\\de cinq},{densité\\de mélange},{conforme\\scindée},{conforme\\normalisée},{conforme\\adaptative}},
  xticklabel style={align=center,font=\footnotesize},tick label style={font=\small},label style={font=\small},
  ymin=70,ymax=95,grid=major,grid style={gray!25},enlarge x limits=0.1,legend columns=2,
  legend style={font=\footnotesize,at={(0.5,-0.3)},anchor=north,draw=none,fill=none,
  /tikz/every even column/.append style={column sep=8pt}},table/col sep=comma]
\addplot[fill=omDef!55,draw=omDef] table[x=name,y=before]{figdata/ml/11-probabilistic-models-and-uncertainty/coverage.csv};
\addplot[fill=omThm!55,draw=omThm] table[x=name,y=after]{figdata/ml/11-probabilistic-models-and-uncertainty/coverage.csv};
\draw[black!70,thick,dashed] ([xshift=-20pt]axis cs:boosted quantiles,90) -- ([xshift=20pt]axis cs:conformal adaptive,90);
\legend{avant le changement de régime,après}
\end{axis}
\end{tikzpicture}

\omcaption{Couverture des intervalles à 90~\% sur les jours de test, avant et après le régime de volatilité (tirets~: le niveau nominal de 90~\%).
Les intervalles propres aux trois modèles, et les intervalles conformes autour de la prévision moyenne des arbres boostés~: \omterm{def:ml:probabilistic-models-and-uncertainty:conformal}{prédiction conforme} scindée avec
erreurs brutes, avec erreurs normalisées par l'écart-type de l'ensemble, et adaptative. Données~: \texttt{ml\_uncert.scores} et
\texttt{ml\_uncert.conformal}.}\label{fig:ml:uncert:coverage}
\end{omfigure}

\begin{omfigure}
\begin{tikzpicture}
\begin{axis}[width=12cm,height=5.4cm,xlabel={jour},ylabel={rendement du lendemain (\%)},tick label style={font=\small},
  label style={font=\small},xmin=2000,xmax=2250,ymin=-7,ymax=7,xtick={2000,2050,2100,2150,2200,2250},grid=major,grid style={gray!25},
  xticklabel style={/pgf/number format/1000 sep={}},legend columns=3,legend style={font=\footnotesize,at={(0.5,-0.3)},
  anchor=north,draw=none,fill=none,/tikz/every even column/.append style={column sep=8pt}},table/col sep=comma]
\addplot[only marks,mark=*,mark size=0.9pt,black!70] table[x=day,y=y]{figdata/ml/11-probabilistic-models-and-uncertainty/band.csv};
\addplot[omDef,thick] table[x=day,y=hi]{figdata/ml/11-probabilistic-models-and-uncertainty/band.csv};
\addplot[omThm,thick,dashed] table[x=day,y=ahi]{figdata/ml/11-probabilistic-models-and-uncertainty/band.csv};
\addplot[omDef,thick,forget plot] table[x=day,y=lo]{figdata/ml/11-probabilistic-models-and-uncertainty/band.csv};
\addplot[omThm,thick,dashed,forget plot] table[x=day,y=alo]{figdata/ml/11-probabilistic-models-and-uncertainty/band.csv};
\draw[gray,thick] (axis cs:2100,-7) -- (axis cs:2100,7);
\legend{rendement réalisé,intervalle à 90~\% de l'ensemble,intervalle conforme adaptatif}
\end{axis}
\end{tikzpicture}

\omcaption{Un actif à travers le changement de régime au jour 2\,100 (ligne verticale)~: rendements réalisés du lendemain, intervalle à 90~\% de
l'ensemble, et l'intervalle conforme adaptatif autour de la prévision moyenne des arbres boostés. Les deux s'élargissent après la rupture~; la bande conforme
s'élargit davantage à mesure que ses erreurs s'accumulent. Données~: \texttt{ml\_uncert.conformal}.}\label{fig:ml:uncert:band}
\end{omfigure}

\section{Utiliser l'incertitude dans le dimensionnement}

Pour des paris indépendants, le poids qui maximise le ratio de Sharpe de leur somme est proportionnel au rendement espéré sur
la variance~; le dimensionnement de Kelly (livre~2, chapitre~29) a la même forme. La variance est là où la distribution du modèle
montre toute son utilité. Le test de dimensionnement du chapitre doit durer plus longtemps que ses trois années de test, car les ratios de Sharpe estimés sur trois
années diffèrent davantage que l'effet~: sur des séries de vingt ans portant sur vingt actifs (quatre graines, évaluées après les sept premières
années), un portefeuille quotidien pondéré par la seule lecture de la dérive obtient un ratio de Sharpe annualisé de 1.72 en moyenne~;
diviser la lecture par une prévision de variance la porte à 1.90 avec la variance EWMA, 1.91 avec la variance prédite par le réseau
gaussien et 1.93 avec la variance vraie. Avec la vraie moyenne, le gain est plus grand (de 2.94 à 3.31)~: plus la
prévision moyenne est bonne, plus la variance compte.

\begin{method}[Prévoir et utiliser une distribution]\label{met:ml:uncert:method}
\begin{enumerate}
\item Entraîner avec une \omterm{def:ml:probabilistic-models-and-uncertainty:scoring}{règle de score propre} (log-vraisemblance négative, \omterm{def:ml:probabilistic-models-and-uncertainty:crps}{pertes pinball}) et évaluer avec une autre (CRPS) et avec
la couverture, sur des périodes hors échantillon, dont une période de stress.
\item Utiliser un ensemble pour séparer l'\omterm{def:ml:probabilistic-models-and-uncertainty:aleatoric}{incertitude épistémique} de l'\omterm{def:ml:probabilistic-models-and-uncertainty:aleatoric}{incertitude aléatoire}~; rapporter les deux.
\item Encadrer la prévision par des intervalles conformes, normalisés par l'échelle prédite, et les mettre à jour en ligne.
\item Dimensionner par la moyenne sur la variance, et tester la règle sur des échantillons assez longs pour que les ratios de Sharpe se distinguent.
\end{enumerate}
\end{method}

\section{Tutoriel~: dimensionner par le doute}\label{tut:ml:probabilistic-models-and-uncertainty:tut}

\begin{tutorial}
\textbf{Objectif.} Ajuster des modèles de \omterm{def:ml:probabilistic-models-and-uncertainty:quantile}{régression quantile}, gaussiens et de mélange au rendement du lendemain, les évaluer avant et après un
régime de volatilité, construire des intervalles conformes scindés et adaptatifs, et dimensionner les positions par la variance prédite.
\textbf{État final~:} le \cref{tab:ml:uncert:scores}, les \cref{fig:ml:uncert:coverage,fig:ml:uncert:band}, les chiffres du dimensionnement.
\begin{tutsteps}
\item \textbf{Scores}~: pinball, CRPS sur une grille de quantiles, couverture.
\omcode{code/firm/uncert/firm_uncert.py}{32}{50}{La perte pinball, le CRPS et la couverture.}\label{lst:ml:uncert:scores}
\item \textbf{Conforme}~: le seuil scindé et le niveau adaptatif en ligne.
\omcode{code/firm/uncert/firm_uncert.py}{164}{185}{Prédiction conforme scindée et adaptative.}\label{lst:ml:uncert:conformal}
\item \textbf{Exécuter} \texttt{ml\_uncert.scores()}, \texttt{conformal()}, \texttt{sizing()} et \texttt{fig\_uncert.py}.
\end{tutsteps}
\textbf{Pour aller plus loin.} Porter $\gamma$ à 0.02 et observer la couverture adaptative se rétablir plus vite et osciller davantage~; dimensionner par
la moyenne sur l'écart-type plutôt que sur la variance.
\end{tutorial}

\section{Construction~: distributions prédictives et intervalles}\label{bld:ml:probabilistic-models-and-uncertainty:uncert}

\begin{build}
\bfield{Rôle} Chaque prévision que la société utilise s'accompagne d'une distribution, d'un score pour l'évaluer et d'un intervalle qui résiste à l'épreuve.
\bfield{Interface} \texttt{pinball}, \texttt{crps\_gaussian}, \texttt{crps\_quantiles}, \texttt{coverage}~;
\texttt{GaussNet}, \texttt{MDN}, \texttt{fit\_nll}, \texttt{predict\_gauss}, \texttt{predict\_mdn}, \texttt{mdn\_quantile}~;
\texttt{ensemble\_split}~; \texttt{split\_conformal(scores, alpha)},
\texttt{adaptive\_conformal(scores, alpha, gamma, window, start)}.
\bfield{Règles} Données de calibration disjointes des données d'entraînement et antérieures au test~; les seuils adaptatifs n'utilisent que les
scores passés~; chaque réseau est déterministe sur un seul thread.
\bfield{Tests d'acceptation} \texttt{code/firm/uncert/tests/}~: la \omterm{def:ml:probabilistic-models-and-uncertainty:crps}{perte pinball} est minimisée au quantile empirique~; le
CRPS gaussien coïncide avec l'approximation sur grille de quantiles~; la \omterm{def:ml:probabilistic-models-and-uncertainty:conformal}{prédiction conforme} scindée couvre au moins $1 - \alpha$ sur des données échangeables
sur de nombreux tirages~; la \omterm{def:ml:probabilistic-models-and-uncertainty:conformal}{prédiction conforme} adaptative rétablit la couverture après un changement d'échelle injecté~; le réseau gaussien et le
modèle de mélange retrouvent une distribution hétéroscédastique et une distribution bimodale injectées.
\bfield{Pour aller plus loin} \omterm{def:ml:probabilistic-models-and-uncertainty:quantile}{Régression quantile} conformalisée~; ensembles conformes pour les classifieurs (étiquettes du chapitre~2)~; une règle de dimensionnement
qui réduit la prévision de variance.
\end{build}

\begin{omsources}
\item T.~Gneiting et A.~E.~Raftery, ``Strictly proper scoring rules, prediction, and estimation'', \emph{Journal of the
American Statistical Association} 102(477), 2007.
\item R.~Koenker et G.~Bassett, ``Regression quantiles'', \emph{Econometrica} 46(1), 1978.
\item C.~M.~Bishop, ``Mixture density networks'', Aston University technical report NCRG/94/004, 1994.
\item V.~Vovk, A.~Gammerman et G.~Shafer, \emph{Algorithmic Learning in a Random World}, Springer, 2005~; 2\up{e} éd., 2022.
\item I.~Gibbs et E.~Candès, ``Adaptive conformal inference under distribution shift'', \emph{NeurIPS}, 2021.
\item A.~Kendall et Y.~Gal, ``What uncertainties do we need in Bayesian deep learning for computer vision?'',
\emph{NeurIPS}, 2017.
\end{omsources}

\section{Exercices}

\begin{exercise}[$\star$]\label{exo:ml:probabilistic-models-and-uncertainty:1}
Une prévision de quantile à 95~\% vaut 2.0~\%~; le rendement est de 3.0~\%. Calculer la \omterm{def:ml:probabilistic-models-and-uncertainty:crps}{perte pinball}. Et si le rendement vaut 1.0~\%~?
\end{exercise}

\begin{exercise}[$\star$]\label{exo:ml:probabilistic-models-and-uncertainty:2}
Avec 250 scores de calibration et $\alpha = 0.1$, quelle statistique d'ordre donne le seuil de \omterm{def:ml:probabilistic-models-and-uncertainty:conformal}{prédiction conforme} scindée~? Et avec 5\,040~?
\end{exercise}

\begin{exercise}[$\star$]\label{exo:ml:probabilistic-models-and-uncertainty:3}
Deux paris ont des rendements espérés de 5 points de base et des écarts types de 80 et 250 points de base. Dans quelle proportion
le dimensionnement moyenne-variance doit-il les détenir~?
\end{exercise}

\begin{exercise}[$\star\star$]\label{exo:ml:probabilistic-models-and-uncertainty:4}
Pourquoi l'intervalle brut de \omterm{def:ml:probabilistic-models-and-uncertainty:conformal}{prédiction conforme} scindée perd-il autant de couverture après la rupture, et l'intervalle normalisé moins~?
\end{exercise}

\begin{exercise}[$\star\star$]\label{exo:ml:probabilistic-models-and-uncertainty:5}
La part épistémique de la variance prédictive vaut 0.09~\%. Que dit-elle de l'intérêt de davantage de données pour cette
prévision, et où s'attendrait-on à ce qu'elle soit plus grande~?
\end{exercise}

\begin{exercise}[$\star\star$]\label{exo:ml:probabilistic-models-and-uncertainty:6}
\emph{Trouver la faille.} «~Le dimensionnement par la variance prédite a fait passer notre ratio de Sharpe de 2.1 à 1.8 sur notre test de
trois ans, donc la mise à l'échelle par la variance ne fonctionne pas pour ce signal.~»
\end{exercise}

\begin{exercise}[$\star\star\star$]\label{exo:ml:probabilistic-models-and-uncertainty:7}
\emph{Programmation.} Exécuter les intervalles de \omterm{def:ml:probabilistic-models-and-uncertainty:conformal}{prédiction conforme} adaptative de \texttt{ml\_uncert.conformal} avec $\gamma = 0.02$. Rapporter la
couverture après la rupture et expliquer le compromis.
\end{exercise}

\begin{exercise}[$\star\star\star$]\label{exo:ml:probabilistic-models-and-uncertainty:8}
Montrer que le CRPS vaut le double de l'intégrale en $\alpha$ de la \omterm{def:ml:probabilistic-models-and-uncertainty:crps}{perte pinball} du $\alpha$-quantile de la prévision.
\end{exercise}

\section{Problème~: dimensionner par le doute}

\begin{problem}[{Problème du week-end --- distributions à travers un changement de régime}]\label{pb:ml:probabilistic-models-and-uncertainty:1}
Les vingt actifs du chapitre, le régime de volatilité à partir du jour 2\,100, et les quatre manières de prévoir le jour suivant.

\medskip\noindent\textbf{Partie I --- Scores.}
\begin{enumerate}
\item Pourquoi un score doit-il être propre, et lesquels des scores du chapitre le sont-ils~?
\item Quels scores obtiennent les modèles (CRPS et couverture) avant la rupture~?
\item Et après~?
\item Pourquoi les CRPS sont-ils si proches d'un modèle à l'autre et de la vérité~?
\end{enumerate}

\medskip\noindent\textbf{Partie II --- Incertitude.}
\begin{enumerate}[resume]
\item Quelle part de la variance prédictive est épistémique~?
\item Qu'apporte un \omterm{def:ml:neural-networks-for-noisy-tabular-data:ensemble}{ensemble profond} par rapport à un réseau gaussien unique ici~?
\item Que modélise le \omterm{def:ml:probabilistic-models-and-uncertainty:quantile}{réseau de densité de mélange} que le réseau gaussien ne modélise pas~?
\item Où s'attendrait-on à ce que la part épistémique compte~?
\end{enumerate}

\medskip\noindent\textbf{Partie III --- \omterm{def:ml:probabilistic-models-and-uncertainty:conformal}{Prédiction conforme}.}
\begin{enumerate}[resume]
\item Que couvrent les intervalles de \omterm{def:ml:probabilistic-models-and-uncertainty:conformal}{prédiction conforme} scindée avant et après la rupture, bruts et normalisés~?
\item Que couvre la \omterm{def:ml:probabilistic-models-and-uncertainty:conformal}{prédiction conforme} adaptative, et à quoi renonce-t-elle~?
\item Quelle hypothèse la rupture viole-t-elle~?
\item Comment choisir $\gamma$~?
\end{enumerate}

\medskip\noindent\textbf{Partie IV --- Dimensionnement et verdict.}
\begin{enumerate}[resume]
\item Pourquoi le test de dimensionnement porte-t-il sur des séries de vingt ans~?
\item Quels ratios de Sharpe les quatre règles obtiennent-elles avec la lecture de la dérive, et avec la vraie moyenne~?
\item Pourquoi la variance compte-t-elle davantage avec une meilleure prévision de la moyenne~?
\item Énoncer le \emph{résultat nommé}~: le gain de Sharpe du dimensionnement par la variance par rapport au dimensionnement par la seule moyenne, et la couverture conforme
face à son niveau nominal de 90~\% après le changement de régime.
\item Que doit montrer un rapport de risque à côté d'une prévision ponctuelle~?
\item Comment suivre les intervalles d'un modèle en production (chapitre~27)~?
\item Que changerait une cible à queues épaisses avec un modèle gaussien~?
\item En une phrase~: qu'est-ce qu'une prévision de rendement~?
\end{enumerate}
\end{problem}

\section{Questions d'entretien}

\begin{interviewq}[$\star$ \iqroles{chercheur, ingénieur ML}]\label{iq:ml:probabilistic-models-and-uncertainty:1}
Qu'est-ce qu'une \omterm{def:ml:probabilistic-models-and-uncertainty:scoring}{règle de score propre}, et pourquoi devriez-vous entraîner et évaluer les prévisions avec une telle règle~?
\end{interviewq}

\begin{interviewq}[$\star\star$ \iqroles{ingénieur ML}]\label{iq:ml:probabilistic-models-and-uncertainty:2}
Comment obtiendriez-vous des intervalles de prédiction à partir d'un modèle à \omterm{def:ml:trees-and-boosting:boosting}{gradient boosting}~?
\end{interviewq}

\begin{interviewq}[$\star\star$ \iqroles{chercheur}]\label{iq:ml:probabilistic-models-and-uncertainty:3}
Expliquez la \omterm{def:ml:probabilistic-models-and-uncertainty:conformal}{prédiction conforme} scindée et sa garantie. Quand la garantie échoue-t-elle en trading~?
\end{interviewq}

\begin{interviewq}[$\star\star$ \iqroles{chercheur, trader}]\label{iq:ml:probabilistic-models-and-uncertainty:4}
Votre modèle prévoit à la fois la moyenne et la variance du rendement de demain. Comment dimensionneriez-vous les positions, et comment
testeriez-vous la règle~?
\end{interviewq}

\begin{interviewq}[$\star\star$ \iqroles{ingénieur ML, chercheur}]\label{iq:ml:probabilistic-models-and-uncertainty:5}
Quelle est la différence entre \omterm{def:ml:probabilistic-models-and-uncertainty:aleatoric}{incertitude aléatoire} et \omterm{def:ml:probabilistic-models-and-uncertainty:aleatoric}{incertitude épistémique}, et comment estimez-vous chacune~?
\end{interviewq}

\begin{interviewq}[$\star\star\star$ \iqroles{chercheur}]\label{iq:ml:probabilistic-models-and-uncertainty:6}
Montrez que la \omterm{def:ml:probabilistic-models-and-uncertainty:crps}{perte pinball} espérée est minimisée par le vrai quantile.
\end{interviewq}
